\chapter{Limits and next steps}
The final test contains two legitimate emails marked as spam and two spam emails allowed through. The saved chart shows how many mistakes occurred, but the project does not yet examine the actual messages behind them. Reviewing those emails would show whether particular words, writing styles or message lengths caused trouble. It would also help decide whether the classifier should be more cautious about blocking legitimate mail or more aggressive about catching spam.

The data comes from an older public collection. Its legitimate messages come from the ``easy ham'' group, which the corpus notes is relatively easy to distinguish from spam\cite{corpus}. The strong test result may therefore be optimistic for modern or unfamiliar email. The project has no separate collection from another time or source for a broader check. Repeated identical emails were removed before evaluation, but very similar messages may remain in different groups.

There is also a small difference between learning and prediction: the prepared learning messages include words identifying the subject and body, while the text sent for a new prediction does not. The project should prepare both in exactly the same way. Future work could also show the estimated chance of spam, examine different decision levels, study mistakes one message at a time, record prediction speed, and test on newer email. These are next steps, not features claimed as complete.

\chapter{Conclusion}
This personal project follows the full path from collecting email data to showing a classification on a web page. The saved comparison shows that the learned models improve greatly on always choosing the most common group. Logistic Regression and SVM tie in the displayed comparison, and the chosen Logistic Regression model makes four mistakes among 590 messages in the reserved final test. The result is promising for this collection, while the age and limited variety of the data leave an important question: how well would it work on new kinds of email? A broader test and a closer look at mistakes are the most useful next steps.
